\section{MXNet: el florecimiento de los marcos de aprendizaje profundo}

\subsection{De CXXNet a MXNet}

El éxito de AlexNet reavivó el entusiasmo de Chen Tianqi (陈天奇) por el aprendizaje profundo. Él y Xu Bin (许斌) (ahora en NVIDIA) desarrollaron CXXNet, que usaba la técnica de programación con plantillas de C++ (expression template) para generar kernels de GPU en línea. La idea central de esta técnica es: cuando escribes la update rule (por ejemplo, operaciones de multiplicación, resta, etc.), las plantillas de C++ pueden comprimir esas operaciones en un solo kernel y generar código de GPU en línea. CXXNet tenía un buen desempeño y también se usó para competir en concursos de Kaggle relacionados con aprendizaje profundo.

Al mismo tiempo, el campo de los marcos de aprendizaje profundo estaba en una etapa de «florecimiento pleno»:

\begin{itemize}
    \item \textbf{Zhang Zhen (张真) y Minjie (敏捷)} (MSR/NYU): desarrollaron Minerva, cuya filosofía de diseño Python first era muy adelantada para su época
    \item \textbf{Li Mu (李沐)} (CMU): trabajó en Parameter Server (servidor de parámetros), que necesitaba que lo distribuido fuera first class citizen
    \item \textbf{Zhang Chiyuan (张池元)} (MIT): escribió en Julia un marco llamado Moka
    \item \textbf{Lin Min (林敏)} (NUS): desarrolló el marco Purine
    \item \textbf{Yang Qing (杨青)} (UC Berkeley): Caffe ya había tenido su release y era el marco dominante de la época, pero su núcleo era C++
\end{itemize}

Gracias a la oportunidad que ofrecían congresos académicos como NeurIPS, estos estudiantes de doctorado dispersos entre UW, CMU, MIT, NUS y NYU se reunieron y decidieron unir fuerzas para construir un marco unificado------este es el origen de MXNet.

\begin{knowledgebox}{El origen del nombre de MXNet}
MXNet es una mezcla (Mix) de Minerva y CXXNet, y al mismo tiempo representa «mix programming style»------la combinación de grafos de cómputo declarativos (declarative, donde primero se declara el grafo de cómputo y después se ejecuta) y de programación imperativa (imperative, donde se calcula en cuanto se escribe). En ese momento existía un trade-off evidente entre estos dos paradigmas: lo declarativo favorecía la optimización del sistema (permitía optimizaciones globales), mientras que lo imperativo era más amigable para el desarrollador (el debug era más simple). MXNet optó por dar soporte a ambos. Más tarde PyTorch eligió el modo puramente imperativo (eager mode), TensorFlow 1.x eligió el modo puramente declarativo (aunque la versión 2.0 después también migró a eager mode), y MXNet ofreció desde el inicio un modo híbrido.
\end{knowledgebox}

Chen Tianqi también mencionó que entre sus compañeros de laboratorio en UW estaba Joe Redmon (el creador de YOLO), «él ya había hecho YOLO antes de su doctorado». Después de la aparición de AlexNet, muchas personas del laboratorio comenzaron a explorar la dirección del aprendizaje profundo.

\subsection{Filosofía central de diseño}

Las tres grandes innovaciones de MXNet:
\begin{enumerate}
    \item \textbf{Python First}: en la época en que Caffe tenía a C++ como núcleo, MXNet eligió a Python como ciudadano de primera clase.
    \item \textbf{Planificación automática}: a partir de las dependencias de cómputo, implementaba de forma automática el paralelismo entre múltiples GPU, de modo que la comunicación y el cómputo pudieran hacer overlap.
    \item \textbf{Diferenciación automática}: evolucionó de la backpropagation escrita a mano en Caffe hacia una diferenciación automática basada en el grafo de cómputo.
\end{enumerate}

\subsection{La competencia con TensorFlow}

Durante un verano en que hacía una pasantía en Google Brain, Chen Tianqi descubrió que TensorFlow estaba a punto de lanzarse. Él fue uno de los usuarios dogfood de TensorFlow, y al mismo tiempo compartió con el equipo de Brain la experiencia de diseño de MXNet. Pero el equipo no se echó atrás: «un ternero recién nacido no le teme al tigre, nosotros también podíamos compete un poco». Incluso siguieron «avanzando con mucha pasión».

En la línea temporal, el proyecto de MXNet se inició antes del lanzamiento de TensorFlow. Todos se habían reunido en el NeurIPS del año anterior para decidir seguir adelante; TensorFlow tuvo su release a finales de 2,015, y el lanzamiento oficial de MXNet fue aproximadamente en la misma época. Al final, MXNet solo envió un workshop paper, con los autores ordenados alfabéticamente------«sencillamente sentíamos que hacer un proyecto juntos era algo bueno».

\begin{knowledgebox}{La colaboración «pura» entre estudiantes de doctorado}
El presentador preguntó: con tantos doctorandos reunidos y sin que nadie estuviera a cargo, ¿cómo podía work out? Chen Tianqi respondió: «en realidad no había nadie que tuviera la última palabra. Porque lo que todos querían, sencillamente, era hacer el proyecto juntos». Este modelo de colaboración basado en intereses compartidos y sin relaciones jerárquicas tenía, en cambio, sus ventajas dentro de la comunidad académica de código abierto------cada quien contribuía por pasión técnica y no por KPI. Dentro de la colaboración, «hablar de la técnica en términos puramente técnicos es en realidad muy simple: se discute la propuesta directamente y se avanza».
\end{knowledgebox}

\begin{warningbox}{Las lecciones del ocaso de MXNet}
MXNet terminó siendo superado por PyTorch, y Chen Tianqi resumió dos lecciones clave:
\begin{enumerate}
    \item \textbf{La experiencia de usuario tiene prioridad sobre el desempeño}: durante mucho tiempo, MXNet fue el mejor marco en paralelismo entre múltiples GPU (NVIDIA lo usaba para ejecutar MLPerf), pero en experiencia de usuario no estaba a la altura de PyTorch. «A la hora de elegir entre experiencia de usuario y desempeño, pensábamos que había que tener ambos, pero en realidad primero había que elegir la experiencia de usuario.»
    \item \textbf{Falta de personal de ingeniería}: dar soporte a cada nueva generación de GPU exigía una gran inversión de ingeniería, y un equipo de estudiantes de doctorado difícilmente podía sostenerlo. Esto lo llevó directamente a orientarse hacia TVM------\textbf{usar tecnología de compilación para reducir la carga de ingeniería}.
\end{enumerate}
\end{warningbox}

\subsection{El papel de pilar de Li Mu (李沐)}

Chen Tianqi mencionó especialmente que Li Mu (李沐) fue quien hizo «un esfuerzo y un sacrificio enormes» en el proyecto MXNet. «En realidad, al final, quien más impulsó todo esto fue Li Mu.» Después de graduarse, Li Mu se fue a Amazon, impulsó que MXNet se convirtiera en el marco oficial de AWS y desarrolló GluonCV y GluonNLP. Chen Tianqi señaló que GluonNLP «in some sense se parece un poco a Hugging Face»------ofrecía una interfaz unificada para distintos modelos preentrenados. Pero, al final, estos esfuerzos a nivel de ecosistema no lograron cambiar el panorama en el que la comunidad de PyTorch ya se había consolidado.

Un hecho poco conocido es que, durante mucho tiempo, NVIDIA estuvo haciendo contribute a MXNet y usándolo. La razón es que, entre todos los marcos de código abierto, el desempeño de MXNet en paralelismo entre múltiples GPU siempre fue el mejor. Al enviar sus resultados de MLPerf, NVIDIA «cuando de verdad necesitaba obtener el mejor puntaje, siempre usaba MXNet»------y solo fue quedando atrás poco a poco después de que ResNet dejó de ser el protagonista del benchmark.

Chen Tianqi también mencionó que Wang Naiyan (王乃岩) y Hou Xiaodi (侯晓迪), de TuSimple (图森未来), «durante mucho tiempo estuvieron ayudando a contribuir y a dar soporte a MXNet». Tanto el éxito de este proyecto como su ocaso final no se pueden entender sin las contribuciones de muchas personas de la comunidad.

\subsection{Resumen de la sección}

La experiencia con MXNet le enseñó dos cosas a Chen Tianqi: (1) la experiencia de usuario importa más que los indicadores técnicos; (2) el cuello de botella de personal de ingeniería debe resolverse con medios técnicos más fundamentales (un compilador), y no simplemente sumando más gente. Estas dos lecciones dieron origen directamente a TVM.
